\section{Traveling salesperson problem}
\label{sec:tsp}

\subsection{Directed graphs}

Let $G_{n, k}$ denote a directed graph with
$
V(G_{n, k}) = \{ v_{c, r} \colon c \in [n], r \in [k] \}
$
and
$$
E(G_{n, k}) = \{ (v_{c, r_1}, v_{(c \mod n) + 1, r_2}) \mid c \in [n], r_1, r_2 \in [k] \}.
$$


It is easy to see that the pathwidth of $G_{n, k}$ is at most $3k$.
Combining the following lemma with \Cref{lem:raz_spieker} immediately gives \Cref{thm:main_dtsp}. The rest of this section will be dedicated to proving it.

\begin{lemma}\label{lem:main_dtsp}
    For every $k \geq 1$ and $n \geq 3k + 3$, any tropical circuit calculating $DTSP_{G_{n, k}}$ is of size at least $C_k / (2k \ln k)$.
\end{lemma}


\begin{proof}
    Fix $n, k$ and put $G := G_{n, k}$. Let $V_i = \{ v_{i, r} : r \in [k] \}$. Let $G^i = G[V_i \cup V_{(i \mod n) + 1}]$ and let $E_i = E(G^i)$, i.e., $E_i$ contains all edges whose tail belongs to $V_i$.
    For every perfect matching $M$ of $G^i$, we will identify it with a permutation $\rho \in \sym_k$ defined as
    $$
    \rho(s) = t
    \quad \text{iff} \quad
    (v_{i, s}, v_{(i \mod n) + 1, t}) \in M.
    $$
    For every Hamiltonian cycle $H$ of $G$, the set $E(H) \cap E_i$ is a perfect matching in $G^i$. We will say that the sequence of permutations $\rho_1, \dots, \rho_n \in \sym_k$ represents $H$ if $E(H) \cap E_i = \rho_i$ for all $i \in [n]$. Note that such representing set is unique, and moreover, satisfies
    $$
    \rho_i \rho_{i - 1} \dots
    \rho_2 \rho_1
    \rho_n \rho_{n - 1} \dots
    \rho_{i + 2} \rho_{i + 1} \in \bar\sym_k
    $$
    for any $i \in [k]$.
    Conversely, every sequence of permutations satisfying the above represents a unique Hamiltonian cycle of $G$.
    In particular, there are $(k - 1)! \cdot \left(k!\right)^{n - 1}$ such sequences and, hence, Hamiltonian cycles in $G$.
    Additionally, for any set of indices $i_1, \dots, i_{\ell} \in [n]$ of size at most $n - 1$, if we fix $\rho_{i_1}, \dots, \rho_{i_{\ell}}$, then there are exactly $(k - 1)! \cdot (k!)^{n - 1 - \ell}$ ways to fix the rest of the permutations for the above inclusion to hold.
    \newline

    Now, look at the polynomial $DTSP_G$. Let $\bar{E} = \{ (v_{i, 1}, v_{(i \mod n) + 1, 1}) \mid i \in [n] \}$ and let $\bar{X} = \{ x_{s,t} \mid (s, t) \in \bar{E} \}$. That is, the set $\bar{X}$ contains exactly one variable corresponding to an edge in $E_i$ for each $i \in [n]$. Similarly, let $X_i$ denote the set $\{ x_{s,t} \mid  (s, t) \in E_i \}$.

    We will say that a pair of polynomials $g, h$ forms a rectangle $g \cdot h$ in $DTSP_G$ if $g \cdot h \subseteq DTSP_G$, i.e., for every pair of monomials $g' \in g, h' \in h$, the variables of $\sup(g') \cup \sup(h')$ correspond to edges forming a Hamiltonian cycle of $G$.
    Additionally, we will say that such rectangle is balanced if $|\sup(g) \cap \bar{X}|, |\sup(h) \cap \bar{X}| \leq \frac{2}{3} |\bar{X}|$. We would like to prove the following claim.
    \begin{claim}\label{lem:dtsp_rect_bound}
        Let $g \cdot h$ be any balanced rectangle in $DTSP_G$. Then
        $$
        |g \cdot h| \leq (k - 1)! \cdot (k!)^{n - 1} \cdot C_k^{-1} \cdot (2k \ln k).
        $$
    \end{claim}

    First, we finish the proof of the lemma given the claim. Let $\tau$ denote the size of the smallest tropical circuit calculating $DTSP_G$. \Cref{lem:tc_decomposition} says that $DTSP_G$ can be covered by a union of $\tau$ balanced rectangles. Thus,
    $$
        (k - 1)! \cdot (k!)^{n - 1} = |DTSP_G| \leq \tau \cdot (k - 1)! \cdot (k!)^{n - 1} \cdot C_k^{-1} \cdot (2k \ln k),
    $$
    hence
    $
    \tau \geq C_k / (2k \ln k).
    $
    The rest of the section is dedicated to proving the claim.
    \newline

    Fix any balanced rectangle $g \cdot h$. Obviously, supports of $g$ and $h$ are disjoint. We will say that $E_i$ is monochromatic w.r.t. $g \cdot h$ if either $X_i \cap \sup(g)$ or $X_i \cap \sup(h)$ is empty. Let $i_1, \dots, i_\ell$ denote the indices of sets $E_i$ that are not monochromatic. We consider two cases depending on whether $\ell > k$ or not.

    Assume $\ell > k$ and fix a non-monochromatic index $i_j$. Define
    $$
    V^g_{tail} = \{ u \in V_{i_j} \mid \exists_{w \in V_{(i_j \mod n) + 1}} x_{u, w} \in \sup(g) \},
    $$
    $$
    V^g_{head} = \{ w \in V_{(i_j \mod n) + 1} \mid \exists_{u \in V_{i_j}} x_{u, w} \in \sup(g) \},
    $$
    i.e., $V^g_{tail}$ (resp. $V^g_{head}$) denote tails (resp. heads) of all edges in $E_{i_j}$ whose related variables belong to the support of $g$. We define $V^h_{tail}, V^h_{head}$ analogously. It is easy to see that $V^g_{tail} \cap V^h_{tail} = V^g_{head} \cap V^h_{head} = \emptyset$.

    A perfect matching in $G^{i_j}$ induced by any monomial of $g \cdot h$ must match $V^g_{head}$ with $V^g_{tail}$ and $V^h_{head}$ with $V^h_{tail}$, and the number of possible ways to do that is
    $
    |V^g_{head}|! \cdot |V^h_{head}|!
    $, which is at most $(k - 1)!$ as both sets are nonempty and disjoint. Therefore, the number of all cycles corresponding to monomials of $g \cdot h$ is at most
    $$
    ((k - 1)!)^\ell \cdot (k!)^{n - \ell} = \frac{1}{k^\ell} \cdot (k!)^n \leq
    \frac{1}{k^{k + 1}} \cdot (k!)^n = (k - 1)! \cdot (k!)^{n - 1} \cdot \frac{1}{k^k}
    $$
    and $\frac{1}{k^k} \leq C_k^{-1} \leq C_k^{-1} \cdot (2k \ln k)$.
    \newline
    
    Now, assume $\ell \leq k$. Thus, at least $n - k \geq \frac{2}{3}n + 1$ of $E_i$ are monochromatic. Since $g \cdot h$ is balanced, we have two indices $i_g, i_h$ such that $E_{i_g}, E_{i_h}$ are both monochromatic, and both intersections $X_{i_g} \cap \sup(h)$ and $X_{i_h} \cap \sup(g)$ are empty. In particular, the support of every monomial in $g$ corresponds to edges whose intersection with $E_{i_g}$ induces a perfect matching in $G^{i_g}$. The same holds for $h$ and $G^{i_h}$. W.l.o.g. we can assume that $i_g < i_h$. Let $\bar{E} = E(G) - (E_{i_g} \cup E_{i_h})$.

    For a monomial $g' \in g$ (and analogously for $h' \in h$), we define its type as the set
    $$
    \lambda_{g'} := \left\{
        e \in \bar{E} \colon
        x_e \in \sup(g')
    \right\}.
    $$
    Fix an arbitrary pair of types $\lambda_{g^*}, \lambda_{h^*}$ and let $g^*$ (resp. $h^*$) denote the set of all monomials of $g$ (resp. $h$) with that type. The number of such possible pairs is bounded by the number of different projections of a Hamiltonian cycle onto $\bar{E}$, hence is at most $(k!)^{n - 2}$. By the definition, every cycle $H$ corresponding to some monomial of $g^* \cdot h^*$ has the same intersection with $\bar{E}$.

    Let $\rho_i$ for $i \in [n] - \{i_g, i_h\}$ represent the set $(\lambda_{g^*} \cup \lambda_{h^*}) \cap E_i$ which is a perfect matching in $G^i$. Let $\bar\rho_l = \rho_{i_h - 1}\rho_{i_h - 2}\dots\rho_{i_g + 1}$ and $\bar\rho_r = \rho_{i_g-1}\rho_{i_g-2}\dots\rho_1\rho_n\dots\rho_{i_h+1}$.
    Let $P_g$ (resp. $P_h$) denote the set of perfect matchings of $G^{i_g}$ (resp. $G^{i_h})$ induced by the monomials in $g^*$ (resp. $h^*$). Every permutation in the product
    $(\bar\rho_l P_g \bar\rho_r) \cdot P_h$ belongs to $\bar\sym_k$, hence by \Cref{lem:mm_size_bound},
    $$
    k(k - 1)! \cdot C_k^{-1} \cdot (2k \ln k) \geq
    |\bar\rho_l P_g \bar\rho_r| \cdot |P_h| =
    |P_g| \cdot |P_h| = |g^*| \cdot |h^*| \geq |g^* \cdot h^*|.
    $$
    Therefore $|g \cdot h| \leq (k!)^{n - 1} \cdot (k - 1)! \cdot C_k^{-1} \cdot (2k \ln k)$.
    
\end{proof}


\subsection{Undirected graphs}

Let $\bar{G}_{n, k}$ denote an undirected graph with
$
V(\bar{G}_{n, k}) = \{ v_{c,r,i} \colon c \in [n], r \in [k], i \in \{-1, 0, 1\} \}
$ and
\begin{align*}
E(\bar{G}_{n, k}) = \{
    (v_{c, r_1, 1}, v_{(c \mod n) + 1, r_2, -1}) &\mid c \in [n], r_1, r_2 \in [k]
\} \ \cup \\
 \{ (v_{c, r, -1}, v_{c, r, 0}) &\mid c \in [n], r \in [k] \} \ \cup \\
 \{ (v_{c, r, 0}, v_{c, r, 1}) &\mid c \in [n], r \in [k] \}.
\end{align*}
The graph $\bar{G}_{n, k}$ is obtained by performing a textbook reduction on $G_{n, k}$ from directed to undirected version of TSP.

\begin{figure}[h]
    \centering
    \label{fig:reduction}
    \begin{subfigure}{0.45\textwidth}
        \centering        \begin{tikzpicture}
		\node [style=node] (0) at (0, 1) [label=$v_{i,1}$] {};
		\node [style=node] (1) at (0, -1) [label=$v_{i,2}$] {};
		\node [style=none] (2) at (-1.5, 1) {};
		\node [style=none] (3) at (-1.5, 0.25) {};
		\node [style=none] (4) at (-1.5, -0.25) {};
		\node [style=none] (5) at (-1.5, -1) {};
		\node [style=none] (6) at (1.5, -1) {};
		\node [style=none] (7) at (1.5, -0.25) {};
		\node [style=none] (8) at (1.5, 0.25) {};
		\node [style=none] (9) at (1.5, 1) {};
        
		\draw [style=directed] (2.center) to (0);
		\draw [style=directed] (3.center) to (1);
		\draw [style=directed] (4.center) to (0);
		\draw [style=directed] (5.center) to (1);
		\draw [style=directed] (1) to (6.center);
		\draw [style=directed] (1) to (8.center);
		\draw [style=directed] (0) to (7.center);
		\draw [style=directed] (0) to (9.center);
\end{tikzpicture}

        \caption{$i$-th column of $G_{n,2}$}
    \end{subfigure}
    \hfill
    \begin{subfigure}{0.45\textwidth}
        \centering        \begin{tikzpicture}
		\node [style=node] (0) at (0, 1) [label={$v_{i,1,0}$}] {};
		\node [style=node] (1) at (0, -1) [label={$v_{i,2,0}$}] {};
		\node [style=node] (2) at (-1, -1) [label={$v_{i,2,-1}$}] {};
		\node [style=node] (3) at (1, -1) [label={$v_{i,2,1}$}] {};
		\node [style=node] (4) at (1, 1) [label={$v_{i,1,1}$}] {};
		\node [style=node] (5) at (-1, 1) [label={$v_{i,1,-1}$}] {};
		\node [style=none] (6) at (-2.5, 1) {};
		\node [style=none] (7) at (-2.5, 0.25) {};
		\node [style=none] (8) at (-2.5, -0.25) {};
		\node [style=none] (9) at (-2.5, -1) {};
		\node [style=none] (10) at (2.5, -1) {};
		\node [style=none] (11) at (2.5, -0.25) {};
		\node [style=none] (12) at (2.5, 0.25) {};
		\node [style=none] (13) at (2.5, 1) {};
        
		\draw [style=undirected] (6.center) to (5);
		\draw [style=undirected] (7.center) to (2);
		\draw [style=undirected] (8.center) to (5);
		\draw [style=undirected] (9.center) to (2);
		\draw [style=undirected] (5) to (0);
		\draw [style=undirected] (0) to (4);
		\draw [style=undirected] (2) to (1);
		\draw [style=undirected] (1) to (3);
		\draw [style=undirected] (4) to (13.center);
		\draw [style=undirected] (3) to (12.center);
		\draw [style=undirected] (4) to (11.center);
		\draw [style=undirected] (3) to (10.center);
\end{tikzpicture}

        \caption{in $\bar{G}_{n,2}$}
    \end{subfigure}
    \caption{The reduction from $G_{n,k}$ to $\bar{G}_{n,k}$}
\end{figure}


Again, the pathwidth of $\bar{G}_{n, k}$ is at most $3k$.
Combining the following lemma with \Cref{lem:raz_spieker} immediately gives \Cref{thm:main_tsp}.

\begin{lemma}\label{lem:tc_tsp}
    For any $k \geq 1$ and $n > 3k$, any tropical circuit calculating $TSP_{\bar{G}_{n, k}}$ uses at least
    %$2^{\Omega(k \log \log k)}$ gates.
    %$2^{\Omega(k \log k)}$ gates.
    $C_k / (2k \ln k)$
    gates.
\end{lemma}

\begin{proof}
    Let $\bar\Gamma$ be an arbitrary tropical circuit of size $\bar\tau$ computing $TSP_{\bar{G}_{n, k}}$. We will show that there exists a tropical circuit $\Gamma$ of size $\tau \leq \bar\tau$ computing $DTSP_{G_{n, k}}$. In particular, together with \Cref{lem:main_dtsp} this shows that $\bar\tau \geq C_k/(2k \ln k)$.

    Let $\Gamma$ be the tropical circuit obtained by taking $\bar\Gamma$ and performing the following substitutions:
    \begin{itemize}
        \item for every input node labeled with a variable $x_{u_1, u_2}$, where $u_1, u_2 \in V(\bar{G}_{n, k})$, $u_1 = v_{c, r, -1}, u_2 = v_{c, r, 0}$, we replace it with a constant $0$,
        \item for every input node labeled with a variable $x_{u_1, u_2}$, where $u_1, u_2 \in V(\bar{G}_{n, k})$, $u_1 = v_{c, r, 0}, u_2 = v_{c, r, 1}$, we replace it with a constant $0$,
        \item for every input node labeled with a variable $x_{u_1, u_2}$, where $u_1, u_2 \in V(\bar{G}_{n, k})$, $u_1 = v_{c, r_1, 1}, u_2 = v_{(c \mod n) + 1, r_2, -1}$, we replace it with a variable $x_{u_1', u_2'}$ where $u_1', u_2' \in V(G_{n, k})$, $u_1' = v_{c, r_1}, u_2' = v_{(c \mod n) + 1, r_2}$.
    \end{itemize}
    Naturally, the size of $\Gamma$ is at most $\bar\tau$. Let $f$ be some polynomial calculated by $\Gamma$. The reduction works in such a way that there is a bijection between directed Hamiltonian cycles in $G_{n, k}$ and undirected ones in $\bar{G}_{n, k}$, and this bijection directly follows the substitution of edges described in the definition of $\Gamma$. In particular, it can be easily seen that $f \simeq DTSP_{G_{n,k}}$.

\end{proof}